\chapter{曲面：两个参数与控制网格}
\label{ch:surfaces}

\section{曲面表达与局部几何}
\sourcepage{第5章}{2--3}

曲线由一个参数定位点，曲面由两个独立参数定位点。球面既可写成隐式形式
\[
 x^2+y^2+z^2-R^2=0,
\]
也可写成参数形式
\[
 \bs S(u,v)=(R\cos u\cos v,\ R\sin u\cos v,\ R\sin v),
 \quad0\leq u\leq2\pi,\quad-\pi/2\leq v\leq\pi/2.
\]
显式形式 $z=\pm\sqrt{R^2-x^2-y^2}$ 必须分为上下两支。参数表达便于规定曲面片范围、计算切向和生成加工路径，但同一曲面可能需要多片参数域，参数化也可能在极点处退化。

对一般曲面 $\bs S(u,v)$，固定 $v$ 或 $u$ 分别得到两族等参数线。偏导矢
\[
 \bs S_u=\frac{\partial\bs S}{\partial u},\qquad
 \bs S_v=\frac{\partial\bs S}{\partial v}
\]
沿两族参数线的切向。若 $\bs S_u\times\bs S_v\ne\bs0$，则为正则点，切平面由这两个矢量张成，单位法矢为
\[
 \bs n=\frac{\bs S_u\times\bs S_v}{\|\bs S_u\times\bs S_v\|}.
\]
叉乘次序决定法矢方向；实际模型还需根据实体内外统一朝向。法矢可用于显示、偏置和刀具姿态分析，但这些操作必须先检查曲面片是否正则。

\nt{B-rep中的一个“面”通常是\textbf{带裁剪边界的曲面片}。曲面方程提供几何信息，边界曲线与相邻关系提供拓扑信息；仅有一组曲面方程，不足以确定一个封闭实体。}

\Needspace{6cm}
\section{Bézier曲面：张量积构造}
\sourcepage{第5章}{4--6}

\dfn{双次数 $(p,q)$ 的Bézier曲面}{由 $(p+1)(q+1)$ 个控制点 $\bs P_{ij}$ 组成控制网格，曲面为
\[
 \bs S(u,v)=\sum_{i=0}^p\sum_{j=0}^q
  B_i^p(u)B_j^q(v)\bs P_{ij},\qquad (u,v)\in[0,1]^2.
\]
固定 $v$ 时，可先计算 $\bs A_i(v)=\sum_jB_j^q(v)\bs P_{ij}$，再求 $\bs S(u,v)=\sum_iB_i^p(u)\bs A_i(v)$。因此可沿两个方向先后应用曲线算法求值。}

\reviewfig[.54]{ch05-bezier-net.pdf}{双三次Bézier曲面的控制网格与两族曲线。第5章讲义，PDF第4页。}{fig:bezier-surface}

\subsection{角点、边界与凸包}
在 $u=0$ 时，仅 $B_0^p(0)=1$，所以
\[
 \bs S(0,v)=\sum_{j=0}^qB_j^q(v)\bs P_{0j}.
\]
其余三个边界同理，分别由对应最外侧控制行或控制列构成Bézier曲线。角点为
\[
 \bs S(0,0)=\bs P_{00},\quad\bs S(1,0)=\bs P_{p0},\quad
 \bs S(0,1)=\bs P_{0q},\quad\bs S(1,1)=\bs P_{pq}.
\]
曲面一般不经过其余控制点；边界曲线也一般不经过该行中间控制点。因为乘积基函数非负且总和为1，曲面位于全部控制点的凸包内，保持仿射不变性。双三次曲面有 $4\times4=16$ 个控制点，两个方向的次数均为3，不能把它误读成“总次数为6的单参数曲线”。

\subsection{偏导与法向}
对 $u$ 方向求导，得到控制点行差分：
\[
 \bs S_u=p\sum_{i=0}^{p-1}\sum_{j=0}^q
 B_i^{p-1}(u)B_j^q(v)(\bs P_{i+1,j}-\bs P_{ij}).
\]
对 $v$ 方向同理：
\[
 \bs S_v=q\sum_{i=0}^p\sum_{j=0}^{q-1}
 B_i^p(u)B_j^{q-1}(v)(\bs P_{i,j+1}-\bs P_{ij}).
\]
例如角点 $(0,0)$ 处，$\bs S_u=p(\bs P_{10}-\bs P_{00})$、$\bs S_v=q(\bs P_{01}-\bs P_{00})$。若这两个控制边共线或有一个为零，该角点参数化退化，不能直接将零叉积归一化。

\ex{双线性曲面的点与法矢}{取
\[
 \bs P_{00}=(0,0,0),\quad\bs P_{10}=(2,0,0),\quad
 \bs P_{01}=(0,2,0),\quad\bs P_{11}=(2,2,1).
\]
双线性插值得到 $\bs S(u,v)=(2u,2v,uv)$。中心为 $(1,1,1/4)$，偏导为 $\bs S_u=(2,0,v)$、$\bs S_v=(0,2,u)$；在中心，叉积为 $(-1,-1,4)$，故
\[
 \bs n=(-1,-1,4)/\sqrt{18}.
\]
双线性表示每个参数分别一次，含混合项 $uv$，所以四个角点不共面时曲面也不必是平面。}

\subsection{曲面拼接：先接边，再比较切平面}
设曲面 $\bs S$ 的 $u=1$ 边界与曲面 $\bs T$ 的 $u=0$ 边界相接，沿边使用同一个参数 $v$。为直接比较，先将公共边的次数及参数化统一。若控制网格分别为 $\bs P_{ij}$ 与 $\bs Q_{ij}$，则 $C^0$ 的直接条件为
\[
 \bs P_{p,j}=\bs Q_{0,j}\quad(\text{所有边界控制点}).
\]
仅四角重合不能保证整条边重合。若两片在跨边方向的次数分别为 $p,r$，参数跨度相同，则跨边一阶偏导相等的充分条件为
\[
 p(\bs P_{p,j}-\bs P_{p-1,j})
 =r(\bs Q_{1,j}-\bs Q_{0,j})\quad(\text{所有 }j),
\]
连同公共边一致，可得到统一参数下的 $C^1$ 连续。若参数跨度不同，仍须按跨度缩放导数。

\reviewfig[.75]{ch05-surface-join.pdf}{两片Bézier曲面的公共控制边及相邻控制行。第5章讲义，PDF第6页。}{fig:surface-join}

$G^1$ 要求的是\textbf{公共边每一点处的切平面连续}，而不强制相同参数速度。在正则条件下，可以用沿边关系表达为
\[
 \bs T_u(0,v)=\alpha(v)\bs S_u(1,v)+\beta(v)\bs S_v(1,v),
 \qquad \alpha(v)>0.
\]
它允许跨边切矢带有沿边分量 $\beta\bs S_v$。一种易检验的充分构造是令所有对应控制点满足
\[
 \bs Q_{1,j}-\bs Q_{0,j}
 =\lambda(\bs P_{p,j}-\bs P_{p-1,j}),\quad\lambda>0
\]
且\textbf{同一个比例 $\lambda$ 对全部 $j$ 一致}。这样沿边的跨界导矢成固定正比例。讲义用“连接前后的控制点形成直线”描述这一几何直观；若每个控制三点虽然共线，却各自使用任意不同的比例，仅凭共线不能无条件保证整条边上都满足 $G^1$。

\section{B样条曲面与NURBS曲面}
\sourcepage{第5章}{7--9}

\subsection{B样条的张量积}
在 $u,v$ 两方向分别指定次数 $p,q$ 与节点向量 $U,V$，定义
\[
 \bs S(u,v)=\sum_{i=0}^n\sum_{j=0}^m
 N_{i,p}(u)M_{j,q}(v)\bs P_{ij}.
\]
$u$ 方向节点数为 $n+p+2$，$v$ 方向为 $m+q+2$，两组节点的数量及间距可分别设置。对某个参数矩形，至多有 $(p+1)(q+1)$ 个控制点起作用，因而能局部编辑。

\begin{itemize}
 \item 沿两个方向分别具有曲线的非负、单位分割、局部支撑和连续性性质。
 \item $u$ 内部节点重数 $a$ 通常保证跨该节点线 $C^{p-a}$；$v$ 方向对应 $C^{q-b}$。两个方向应分别判断。
 \item 只有两方向都端点夹持时，四个角点才插值角部控制点；夹持边界为相应最外侧控制行或列的B样条曲线。
 \item 若 $n=p,m=q$，且两方向均使用仅含端点的夹持节点，便退化为Bézier曲面。只在一个方向满足时，另一方向仍可保留多段B样条结构。
\end{itemize}
求值可以先沿所有控制行使用de Boor，再沿另一方向求值；节点插入与升阶也可逐行或逐列实施。\textbf{张量积结构要求规则的控制点网格}，不等于输入任意散乱点就自动得到一张曲面。

\subsection{NURBS曲面的权重}
给每个控制点一个权重 $w_{ij}$，定义
\[
 \bs S(u,v)=
 \frac{\displaystyle\sum_{i=0}^n\sum_{j=0}^m
       w_{ij}N_{i,p}(u)M_{j,q}(v)\bs P_{ij}}
      {\displaystyle\sum_{i=0}^n\sum_{j=0}^m
       w_{ij}N_{i,p}(u)M_{j,q}(v)}.
\]
全部权重相同时退化为B样条曲面；通常取正权重，使分母非零并保持凸包性。NURBS能统一表达自由曲面与圆柱、圆锥、球面等二次曲面。其精确圆形截面来源于有理曲线，而不是靠大量三角面片逼近圆。

\subsection{从曲线扫掠生成曲面}
最简单的直线扫掠是将一条NURBS曲线 $\bs C(v)$ 沿固定矢量 $\bs d$ 拉伸：
\[
 \bs S(u,v)=\bs C(v)+u\bs d,\qquad0\leq u\leq1.
\]
若曲线控制点为 $\bs P_j$、权重 $w_j$，可设置两行控制点
\[
 \bs P_{0j}=\bs P_j,\quad\bs P_{1j}=\bs P_j+\bs d,
 \qquad w_{0j}=w_{1j}=w_j.
\]
$u$ 方向取次数1、阶次2和节点 $(0,0,1,1)$，$v$ 方向沿用母线的次数、节点与权重，就得到精确的拉伸曲面。

\reviewfig[.70]{ch05-sweep.pdf}{沿直线移动曲线生成扫掠曲面；两个参数分别描述移动位置与母线位置。第5章讲义，PDF第8页。}{fig:sweep-surface}

更一般的扫掠沿路径移动并可能旋转截面，必须规定截面姿态。\textbf{旋转曲面}则将母线绕轴旋转；圆弧方向使用有理二次表示，才能精确生成圆柱、圆锥等曲面。不要把所有扫掠曲面都套入上述固定矢量拉伸公式。

\Needspace{11cm}
\section{复习速查与自测}
\begin{enumerate}
 \item 双三次Bézier曲面有多少个控制点，能保证通过哪些点？\textit{答：16个；四角控制点。边界是Bézier曲线，一般不经过边界中间控制点。}
 \item 将 $u$ 固定得到什么？\textit{答：一条关于 $v$ 的等参数曲线；不是曲面法线。}
 \item 曲面法矢为何采用两个偏导的叉积？\textit{答：两个偏导张成切平面，其叉积垂直于该平面；正则性要求叉积非零。}
 \item 两曲面的公共边端点重合，是否已有 $C^0$ 连续？\textit{答：不足；整条边界曲线必须重合。}
 \item 两曲面达到 $G^1$ 就一定达到 $C^1$ 吗？\textit{答：不一定；$G^1$ 只要求切平面连续，$C^1$ 还要求同一参数下相应一阶偏导连续。}
 \item 有 $6\times5$ 个控制点的双三次B样条曲面，两方向各需多少节点？\textit{答：$u$ 为 $6+3+1=10$，$v$ 为 $5+3+1=9$。}
 \item 将NURBS曲面所有权重同乘2有何变化？\textit{答：无变化。只改变局部权重一般会改变曲面形状。}
 \item 将一条三次NURBS曲线沿直线拉伸，另一方向的次数是否也为3？\textit{答：不必，直线拉伸方向只需次数1、两个控制点。}
\end{enumerate}

\nt{解曲面题时，先标清 $u,v$ 与控制网格下标的对应，再分别应用曲线结论。涉及边界与连续性时明确是哪个方向、哪条边，以及比较几何形状还是比较参数导数。}
